%% kapa: exam=1 round=30 year=2019 period=2 form=A subject=appraisal_law subject_ko=감정평가관계법규 seq=24 number=24
%% source: Q-Net 감정평가사 기출문제 (cst00309) artlSeq=5205205 file=extracted/1008/5205205/2019년 제30회 감정평가사 1차시험 문제지/2019년 제30회 감정평가사 1차시험 문제지/2019년도 제30회 감정평가사 1차 2교시 A형.hwp
%% conversion: kordoc 4.24.51 -> tools/gen1.py (auto); review=proofread-sonnet; flags=none
\begin{question}{24}
건축법령상 특별건축구역에 관한 설명으로 옳은 것은?
\begin{choices}
  \item 시장ㆍ군수ㆍ구청장은 특별건축구역의 지정을 신청할 수 없다.
  \item 「군사기지 및 군사시설 보호법」에 따른 군사기지 및 군사시설 보호구역은 특례 적용이 필요하다고 인정하는 경우에도 특별건축구역으로 지정될 수 없다.
  \item 시ㆍ도지사는「도시개발법」에 따른 도시개발구역에 대하여 특별건축구역을 지정할 수 있다.
  \item 「주택도시기금법」에 따른 주택도시보증공사가 건축하는 건축물은 특별건축구역에서 특례사항을 적용하여 건축할 수 있는 건축물에 해당된다.
  \item 지정신청기관은 특별건축구역 지정 이후 특별건축구역의 도시ㆍ군관리계획에 관한 사항이 변경되는 경우에는 변경지정을 받지 않아도 된다.
\end{choices}
\end{question}
\answer{3}
